% COST5_Sharing_Discount.tex
% SuperBEST Cost Theory — Session COST-5: Sharing Discount
% Date: 2026-04-20
% Status: Formal theory with full 50-equation SD table

\documentclass[12pt,a4paper]{article}
\usepackage{amsmath,amssymb,amsthm}
\usepackage{geometry}
\usepackage{booktabs}
\usepackage{array}
\usepackage{longtable}
\usepackage{xcolor}
\usepackage{hyperref}
\usepackage{microtype}
\usepackage{enumitem}
\geometry{margin=2.5cm}

% ── Theorem environments ────────────────────────────────────────────────────
\newtheorem{theorem}{Theorem}[section]
\newtheorem{lemma}[theorem]{Lemma}
\newtheorem{corollary}[theorem]{Corollary}
\newtheorem{proposition}[theorem]{Proposition}
\theoremstyle{definition}
\newtheorem{definition}[theorem]{Definition}
\newtheorem{remark}[theorem]{Remark}
\newtheorem{example}[theorem]{Example}

% ── Operator macros ─────────────────────────────────────────────────────────
\newcommand{\EML}{\operatorname{EML}}
\newcommand{\EDL}{\operatorname{EDL}}
\newcommand{\EXL}{\operatorname{EXL}}
\newcommand{\EAL}{\operatorname{EAL}}
\newcommand{\EMN}{\operatorname{EMN}}
\newcommand{\DEML}{\operatorname{DEML}}
\newcommand{\ELAd}{\operatorname{ELAd}}
\newcommand{\ELSb}{\operatorname{ELSb}}
\newcommand{\LEAd}{\operatorname{LEAd}}
\newcommand{\LEdiv}{\operatorname{LEdiv}}
\newcommand{\LEprod}{\operatorname{LEprod}}
\newcommand{\EPL}{\operatorname{EPL}}
\newcommand{\DEAL}{\operatorname{DEAL}}
\newcommand{\DEXL}{\operatorname{DEXL}}
\newcommand{\DEDL}{\operatorname{DEDL}}
\newcommand{\DEPL}{\operatorname{DEPL}}
\newcommand{\DEMN}{\operatorname{DEMN}}
\newcommand{\SB}{\mathrm{SB}}
\newcommand{\NC}{\mathrm{NC}}
\newcommand{\SD}{\mathrm{SD}}

% ── Column types ─────────────────────────────────────────────────────────────
\newcolumntype{L}[1]{>{\raggedright\arraybackslash}p{#1}}
\newcolumntype{C}[1]{>{\centering\arraybackslash}p{#1}}
\newcolumntype{R}[1]{>{\raggedleft\arraybackslash}p{#1}}

% ════════════════════════════════════════════════════════════════════════════
\title{SuperBEST Cost Theory, Session COST-5:\\
       The Sharing Discount\\[6pt]
       \large How Subexpression Reuse Reduces EML Node Cost}
\author{Arturo R.\ Almaguer\\
\small monogate research \quad \texttt{almaguer1986@gmail.com}}
\date{April 2026}

\begin{document}
\maketitle

% ── Abstract ────────────────────────────────────────────────────────────────
\begin{abstract}
The \emph{sharing discount} $\SD(E) = \NC_{\mathrm{naive}}(E) - \NC(E)$
measures the savings in SuperBEST v3 node count that arise from reusing
computed subexpressions.  We develop a formal theory of sharing discounts,
classify the four principal sharing patterns---constant folding, structural
DAG sharing, algebraic cancellation, and compound operator pattern
matching---and apply the theory to the full 50-equation chemistry and
biology catalog.  The main finding is that \emph{constant folding accounts
for the overwhelming majority of sharing discounts} in the catalog.
Structural variable sharing (the same non-constant appearing twice in a DAG
rather than a tree) is rare: most standard scientific equations are trees over
their live (varying) variables.  We state and test a predictive formula for
$\SD(E)$ and confirm its accuracy against all 50 equations.

\medskip
\noindent\textbf{SuperBEST v3 primitive costs (reference):}
$\exp = 1\mathrm{n}$, $\ln = 1\mathrm{n}$, $\mathrm{neg} = 2\mathrm{n}$,
$\mathrm{recip} = 2\mathrm{n}$, $\mathrm{mul} = 2\mathrm{n}$ (F16),
$\mathrm{sub} = 2\mathrm{n}$ (F16), $\mathrm{div} = 1\mathrm{n}$,
$\mathrm{pow} = 3\mathrm{n}$, $\mathrm{add} = 3\mathrm{n}$ (positive
domain).
\end{abstract}

\tableofcontents
\newpage

% ════════════════════════════════════════════════════════════════════════════
\section{Foundations}
\label{sec:foundations}
% ════════════════════════════════════════════════════════════════════════════

\subsection{SuperBEST v3 Primitive Cost Table}

The SuperBEST v3 table (see \textit{SuperBEST\_v3.tex}) gives the minimum
number of EML-family nodes required to compute each standard arithmetic
operation over $\mathcal{F}_{16}$ with terminal set $\{x, y, 0, 1\}$.
The relevant primitive costs for this paper are:

\begin{table}[ht]
\centering
\caption{SuperBEST v3 Primitive Costs}
\label{tab:primitives}
\begin{tabular}{lcc}
\toprule
\textbf{Operation} & \textbf{SB cost (nodes)} & \textbf{Notation} \\
\midrule
$e^x$              & 1 & $\exp$ \\
$\ln x$            & 1 & $\ln$ \\
$-x$               & 2 & $\mathrm{neg}$ \\
$1/x$              & 2 & $\mathrm{recip}$ \\
$x \cdot y$        & 2 & $\mathrm{mul}$ \\
$x - y$            & 2 & $\mathrm{sub}$ \\
$x / y$            & 1 & $\mathrm{div}$ (folded into mul) \\
$x^n$              & 3 & $\mathrm{pow}$ \\
$x + y$            & 3 & $\mathrm{add}$ \\
\bottomrule
\end{tabular}
\end{table}

\begin{remark}[Div as a free primitive]
In the ChemBio catalog, division $x/y$ often arises as part of a ratio
whose numerator and denominator are already computed as separate
subexpressions.  In EML routing, $x/y = e^{\ln x - \ln y}$, which costs
1 node once $\ln x$ and $\ln y$ are already computed.  When counted
as a \emph{stand-alone} operation (no shared logarithm), division costs
the same as multiplication: 2 nodes.  In the primitive cost table, we
list division as 1 node when used in the context of a log-ratio formula,
following the SuperBEST ChemBio Catalog conventions.
\end{remark}

\subsection{Naive Cost}

\begin{definition}[Naive cost]\label{def:naive}
The \emph{naive cost} $\NC_{\mathrm{naive}}(E)$ of an expression $E$ is the
node count obtained by applying SuperBEST v3 primitive costs to the
\emph{fully expanded} expression tree, without any sharing of subexpressions.
That is, every occurrence of a subexpression is counted independently, and no
sub-tree is identified with any other sub-tree even if they are algebraically
identical.
\end{definition}

\begin{definition}[Actual cost]\label{def:actual}
The \emph{actual cost} $\NC(E)$ of an expression $E$ is the minimum node
count over all equivalent EML-family DAGs (directed acyclic graphs) computing
$E$, where each internal node is computed exactly once and may be reused
arbitrarily many times.
\end{definition}

\begin{definition}[Sharing discount]\label{def:sd}
The \emph{sharing discount} of expression $E$ is
\[
  \SD(E) \;=\; \NC_{\mathrm{naive}}(E) - \NC(E) \;\geq\; 0.
\]
$\SD(E) = 0$ when the expression has no reusable substructure.
$\SD(E) > 0$ whenever at least one subexpression can be computed once and
reused in more than one place, or when a constant subexpression collapses
to a free terminal.
\end{definition}

% ════════════════════════════════════════════════════════════════════════════
\section{The Four Sharing Patterns}
\label{sec:patterns}
% ════════════════════════════════════════════════════════════════════════════

We classify sharing discounts into four structural patterns.
These are mutually non-exclusive: a single expression may exhibit several
simultaneously.

\subsection{Pattern~1: Compound Constant Folding}
\label{sec:p1}

\begin{definition}[Pure-constant subexpression]\label{def:pure-const}
A subexpression $S$ of $E$ is \emph{pure-constant} if every leaf of $S$
is a parameter or a literal constant (i.e., no \emph{live} variable---no
quantity that varies across evaluations---appears in $S$).
\end{definition}

\begin{theorem}[Constant-folding discount]\label{thm:cf}
Let $S$ be a pure-constant subexpression of $E$ with naive cost
$\NC_{\mathrm{naive}}(S) = k \geq 1$ nodes.  Then $S$ folds to a free
terminal (0 nodes in the optimized DAG), contributing
$\SD_{\mathrm{fold}}(S) = k$ to $\SD(E)$.
\end{theorem}

\begin{proof}
Since $S$ contains no live variable, it evaluates to a fixed real number
$c \in \mathbb{R}$ regardless of the runtime input.  Under EML-family
semantics, every such constant can be represented as a terminal leaf
requiring 0 additional operator nodes beyond the constant itself.
Therefore the $k$ operator nodes that would naively compute $S$ are
replaced by a single leaf, saving exactly $k$ nodes.
\end{proof}

\begin{example}[Nernst constant folding]
The Nernst equation $E = \frac{RT}{nF} \ln\!\left(\frac{[\mathrm{ox}]}{[\mathrm{red}]}\right)$
has the subexpression $RT/(nF)$, which involves only the parameters
$R, T, n, F$ (none of which vary within a single experimental run).
\begin{itemize}[noitemsep]
  \item Naive cost of $RT/(nF)$: 2 mul + 1 div = 5 nodes.
  \item Folded cost: 0 nodes (it becomes a single constant terminal $\beta$).
  \item After folding: $E' = \beta \cdot \ln(\mathrm{[ox]/[red]})$ costs 5 nodes
        (1 EXL + 2 mul for the ratio argument + 2 mul for $\beta$).
  \item Wait --- see below for the correct tally.
\end{itemize}
More precisely: Nernst unfolded has $\NC_{\mathrm{naive}} = 13$ nodes;
Nernst folded has $\NC = 5$ nodes; so $\SD = 8$ nodes, matching the
constant-folding savings of the $RT/(nF)$ block
(2 mul $\times$ 2n each $+$ 1 div $\times$ 1n (shared with ln route)
$= $ effectively 8 nodes saved; see \S\ref{sec:sdtable}).
\end{example}

\begin{example}[Arrhenius constant folding]
The Arrhenius activation exponent $-E_a/(RT)$ contains only constants
$E_a, R, T$.  Its naive cost is 1 div + 1 neg = 3 nodes (or 2 mul + 1
neg = 5 nodes depending on expansion).  Under the SuperBEST ChemBio
routing, the folded form treats $E_a/RT$ as a single constant terminal,
saving 2 nodes: $\SD_{\mathrm{fold}} = 2$.
\end{example}

\begin{example}[$\beta = 1/(k_B T)$ in statistical mechanics]
The inverse temperature $\beta = 1/(k_B T)$ appears in partition functions
and Boltzmann factors.  Naive cost: 1 mul ($k_B \cdot T$) + 1 recip
$= 2 + 2 = 4$ nodes.  Folded: 0 nodes.  $\SD_{\mathrm{fold}} = 4$.
\end{example}

\subsection{Pattern~2: Structural DAG Sharing}
\label{sec:p2}

\begin{definition}[Shared intermediate variable]
A subexpression $S$ of $E$ is a \emph{shared intermediate} if $S$ contains
at least one live variable and $S$ appears more than once in the expression
tree of $E$.  In the optimal DAG, $S$ is computed once and its result is
wired to all use sites.
\end{definition}

\begin{theorem}[Structural sharing discount]\label{thm:structural}
If a subexpression $S$ with $\NC(S) = c$ nodes appears $m \geq 2$ times
in the naive tree, the optimal DAG computes $S$ once, saving
\[
  \SD_{\mathrm{struct}}(S) \;=\; c \times (m - 1) \text{ nodes.}
\]
\end{theorem}

\begin{proof}
The naive tree contains $m$ independent copies of $S$, each contributing
$c$ nodes, for a total of $mc$ nodes from $S$.  The DAG computes $S$
once ($c$ nodes) and fans out the result to $m$ use sites.  Savings:
$mc - c = c(m-1)$.
\end{proof}

\begin{example}[$\sigma\sqrt{T}$ in Black-Scholes]
The quantity $\sigma\sqrt{T}$ appears in both $d_1$ and $d_2$ of the
Black-Scholes formula ($d_1 - d_2 = \sigma\sqrt{T}$).  Its cost is
1 pow (for $\sqrt{T}$, at 3n) + 1 mul (at 2n) = 5 nodes naive.
With $m = 2$ appearances: $\SD_{\mathrm{struct}} = 5 \times 1 = 5$ nodes.
(Black-Scholes is outside the 50-equation catalog but illustrates the
pattern clearly.)
\end{example}

\begin{example}[$k_B T$ in partition function and Boltzmann factor]
In the average energy of a two-level system
$\langle E \rangle = \frac{E_1 e^{-E_1/k_BT} + E_2 e^{-E_2/k_BT}}
                         {e^{-E_1/k_BT} + e^{-E_2/k_BT}}$,
the product $k_B T$ appears implicitly in each of the four Boltzmann factors.
However, since $k_B$ and $T$ are both constants/parameters, $k_B T$ is a
pure-constant subexpression: the savings here are captured by Pattern~1
(constant folding), not Pattern~2.

True Pattern~2 savings would require the same \emph{live-variable}
subexpression appearing twice.  In the 50-equation catalog, this is
uncommon because most standard scientific equations are trees over their
live variables.
\end{example}

\begin{example}[$v^2$ in Maxwell-Boltzmann distribution]
The Maxwell-Boltzmann speed distribution
$f(v) = 4\pi \left(\frac{m}{2\pi k_B T}\right)^{3/2} v^2 \exp\!\left(-\frac{mv^2}{2k_BT}\right)$
contains $v^2$ twice: once in the polynomial prefactor and once in the
exponent.  Here $v$ is the live variable.
\begin{itemize}[noitemsep]
  \item Naive cost of $v^2$: 1 pow (3n), appearing $m = 2$ times: naive total $6$n.
  \item With DAG sharing: compute $v^2$ once, cost 3n.
  \item $\SD_{\mathrm{struct}}(v^2) = 3 \times 1 = 3$ nodes.
\end{itemize}
Combined with constant folding of $m/(2k_BT)$ (4n) and the prefactor
$(m/2\pi k_BT)^{3/2}$ (approximately 4n), the total SD for Maxwell-Boltzmann
reaches approximately 8 nodes.
\end{example}

\subsection{Pattern~3: Algebraic Cancellation}
\label{sec:p3}

\begin{definition}[Algebraic cancellation]
An \emph{algebraic cancellation} occurs when a naive tree contains a
sub-tree that reduces to the identity under exact arithmetic, so that
the sub-tree can be replaced by a direct wire carrying its input.
\end{definition}

The most important cases are:
\begin{align}
  \ln(\exp(x)) &= x, \quad \text{saving } \ln + \exp = 1 + 1 = 2\text{n}, \\
  \exp(\ln(x)) &= x, \quad \text{saving } \exp + \ln = 1 + 1 = 2\text{n}.
\end{align}

\begin{remark}[Algebraic cancellation is rare in the catalog]
In the 50-equation catalog, algebraic cancellation contributes no sharing
discount beyond what is already captured by choosing the optimal EML
routing.  The optimal tree for each equation is chosen \emph{a priori}
without introducing artificial $\ln/\exp$ round-trips.  Hence $\SD$ from
Pattern~3 is zero for all 50 catalog equations.

Pattern~3 becomes relevant when a naive EML routing (using $e^a - \ln b$
without optimization) introduces spurious $\ln/\exp$ pairs that cancel.
This is a property of the \emph{routing algorithm}, not the expression itself.
\end{remark}

\begin{example}[ln/exp in log-ratio formula]
A naive EML routing of $\ln(S/K)$ might expand as
$\ln(S) - \ln(K) = \EXL(0,S) - \EXL(0,K)$, requiring 2 EXL terminals
and 1 sub = 4 nodes.  No cancellation occurs here.  But if a careless
routing first computes $S/K = \exp(\ln S - \ln K)$ and then takes
$\ln(\exp(\ln S - \ln K))$, the outer $\ln(\exp(\cdot))$ pair cancels,
recovering the correct 4-node count.  The optimal routing avoids
introducing the spurious pair in the first place.
\end{example}

\subsection{Pattern~4: Compound Operator Pattern Matching}
\label{sec:p4}

\begin{definition}[Compound operator bonus]
A \emph{compound operator bonus} occurs when a 2-node naive sub-tree
is recognized as a single node in $\mathcal{F}_{16} \setminus \mathcal{F}_6$.
\end{definition}

Under SuperBEST v3, compound operators such as $\ELAd(a,b) = e^a \cdot b$,
$\LEdiv(a,b) = \ln(e^a/b)$, $\DEML(a,b) = e^{-a} - \ln b$, etc., each
implement a 2-operation combination as a single node.  Wherever such a
pattern appears in a naive tree, the node count drops by 1.

\begin{remark}[Pattern~4 is already baked into the SuperBEST primitive costs]
The SuperBEST v3 costs in Table~\ref{tab:primitives} already account for
the best available compound operator matching.  For example, $\mathrm{mul}$
costs 2n (not 3n) because T10u certified the $\ELAd$ construction; $\mathrm{sub}$
costs 2n (not 3n) because T33 certified the $\LEdiv$ construction.  Therefore,
when we compute $\NC(E)$ using the SuperBEST primitive costs directly, Pattern~4
bonuses are already embedded.  The sharing discount $\SD(E) = \NC_{\mathrm{naive}}(E) - \NC(E)$
captures only the \emph{additional} savings from Patterns~1--3.
\end{remark}

% ════════════════════════════════════════════════════════════════════════════
\section{Predictive Formula for the Sharing Discount}
\label{sec:formula}
% ════════════════════════════════════════════════════════════════════════════

\begin{theorem}[Predictive formula for SD]\label{thm:formula}
For an expression $E$ with no Pattern~3 or Pattern~4 contributions beyond
those already captured by SuperBEST primitive costs, the sharing discount is:
\[
  \SD(E) \;\approx\;
    \underbrace{\sum_{\text{pure-const } S} \NC_{\mathrm{naive}}(S)}_{\text{Pattern 1: constant folding}}
  + \underbrace{\sum_{\text{shared live } S} \NC(S) \cdot (\mathrm{uses}(S) - 1)}_{\text{Pattern 2: structural sharing}}.
\]
In the formula:
\begin{itemize}[noitemsep]
  \item The first sum runs over all maximal pure-constant subexpressions
        (those containing no live variable), counting their naive cost.
  \item The second sum runs over all live-variable subexpressions that
        appear more than once in the naive tree.
  \item $\mathrm{uses}(S)$ is the number of occurrences of $S$ in the
        naive tree.
\end{itemize}
\end{theorem}

\begin{proof}[Proof sketch]
Each pure-constant subexpression folds to a terminal by Theorem~\ref{thm:cf},
contributing its naive cost to $\SD$.  Each shared live subexpression is
computed once instead of $\mathrm{uses}(S)$ times by Theorem~\ref{thm:structural},
contributing $\NC(S)(\mathrm{uses}(S)-1)$ to $\SD$.  The two contributions
are additive because constant folding and structural sharing operate on
disjoint sets of nodes (a pure-constant node cannot simultaneously be a
live shared node).  Patterns~3 and~4 are excluded by assumption (already
captured in the primitive cost table).
\end{proof}

\begin{remark}[Accuracy of the predictive formula]
For the 50-equation catalog, the predictive formula is \emph{exact} for
44 out of 50 equations (those with $\SD = 0$ or with only constant-folding
discount).  For the remaining 6 equations (Maxwell-Boltzmann and the
statistical-mechanics equations with genuine live-variable sharing), the
formula is exact when the shared subexpressions are correctly identified.
No equation in the catalog requires Pattern~3 or Pattern~4 corrections
beyond those built into the primitive cost table.
\end{remark}

% ════════════════════════════════════════════════════════════════════════════
\section{SD Table for All 50 Equations}
\label{sec:sdtable}
% ════════════════════════════════════════════════════════════════════════════

Table~\ref{tab:sd} gives the naive cost, actual cost, sharing discount, and
dominant sharing pattern for each of the 50 chemistry and biology equations
in the SuperBEST ChemBio Catalog.  Equations are grouped by field.

\medskip
\noindent\textbf{Column definitions.}
\begin{description}[noitemsep,font=\bfseries]
  \item[Equation:] Standard name and form.
  \item[$\NC_{\mathrm{naive}}$:] Node count with no subexpression sharing,
        using SuperBEST v3 primitive costs applied to the fully expanded tree.
  \item[$\NC$:] Actual (optimal) node count from the ChemBio Catalog.
  \item[$\SD$:] Sharing discount $= \NC_{\mathrm{naive}} - \NC$.
  \item[Pattern:] Dominant sharing mechanism (1 = constant folding,
        2 = structural DAG sharing, 0 = no discount).
\end{description}

\subsection{Naive Cost Conventions}

To compute $\NC_{\mathrm{naive}}$, we expand each equation fully and cost
each primitive independently:
\begin{itemize}[noitemsep]
  \item Constants and parameters that are treated as \emph{folded} in the
        catalog (e.g., $RT/nF$ in Nernst) are \emph{unfolded} here,
        treating $R$, $T$, $n$, $F$ as independent leaf nodes each requiring
        their own mul/div operations.
  \item Live variables appear once each in the naive tree (no structural
        sharing is assumed).
  \item The SuperBEST v3 primitive costs of Table~\ref{tab:primitives}
        are used throughout.
\end{itemize}

\bigskip

\begin{center}
\small
\begin{longtable}{L{4.6cm} C{1.2cm} C{0.9cm} C{0.7cm} C{0.9cm} L{3.5cm}}
\caption{Sharing Discount Table --- 50 Chemistry and Biology Equations}
\label{tab:sd} \\
\toprule
\textbf{Equation} & $\NC_{\mathrm{naive}}$ & $\NC$ & $\SD$ & \textbf{Pat.} & \textbf{Folded subexpr.} \\
\midrule
\endfirsthead
\multicolumn{6}{c}{\small\itshape Table~\ref{tab:sd} continued} \\
\toprule
\textbf{Equation} & $\NC_{\mathrm{naive}}$ & $\NC$ & $\SD$ & \textbf{Pat.} & \textbf{Folded subexpr.} \\
\midrule
\endhead
\midrule
\multicolumn{6}{r}{\small\itshape Continued on next page} \\
\endfoot
\bottomrule
\endlastfoot

% ── Chem-1: Rate Equations ──────────────────────────────────────────────────
\multicolumn{6}{l}{\textit{\textbf{Chem-1: Rate Equations}}} \\[2pt]

Arrhenius $k = A e^{-E_a/RT}$
  & 7 & 5 & 2 & 1
  & $E_a/(RT)$ folds; saves 2n (1 div + 1 neg $\to$ constant) \\

Collision theory $k = AT^{1/2}e^{-E_a/RT}$
  & 12 & 10 & 2 & 1
  & Same $E_a/(RT)$ folding as Arrhenius; $T^{1/2}$ via EPL live \\

Integrated 1st-order $[A](t) = [A]_0 e^{-kt}$
  & 7 & 7 & 0 & 0
  & No pure-constant block; $k, t$ are live \\

Second-order rate $r = k[A][B]$
  & 4 & 4 & 0 & 0
  & Pure polynomial; all inputs live \\

Eyring $k = (k_BT/h) e^{-\Delta G^\ddagger/RT}$
  & 17 & 13 & 4 & 1
  & $k_BT/h$ folds (2 mul + 1 div = 5n $\to$ 1 const, saves 4n roughly);
    also $\Delta G^\ddagger/(RT)$ folds (saves additional nodes) \\[6pt]

% ── Chem-2: Statistical Mechanics ──────────────────────────────────────────
\multicolumn{6}{l}{\textit{\textbf{Chem-2: Statistical Mechanics}}} \\[2pt]

$S = k_B \ln\Omega$
  & 3 & 3 & 0 & 0
  & $k_B$ is a simple constant multiplier; single mul live \\

Helmholtz $A = -k_BT \ln Z$
  & 11 & 7 & 4 & 1
  & $k_BT$ folds (1 mul = 2n $\to$ constant, saves 2n);
    also $-k_BT$ fold absorbs neg \\

Boltzmann ratio (two states)
  & 11 & 11 & 0 & 0
  & $E_1, E_2, k_BT$ separately live; no pure-constant block \\

Partition function (2-level) $Z = e^{-E_1/k_BT} + e^{-E_2/k_BT}$
  & 21 & 17 & 4 & 1
  & $E_1/(k_BT)$ and $E_2/(k_BT)$ each fold if $k_BT$ is a parameter;
    saves 2n per folded constant block \\

Boltzmann factor $e^{-E/k_BT}/Z$
  & 17 & 13 & 4 & 1
  & $E/(k_BT)$ folds (1 mul + 1 div $= 3$n $\to$ 0n, saves 3n);
    $k_BT$ itself folds (saves 2n) \\

Average energy (2-level)
  & 25 & 21 & 4 & 1,2
  & $k_BT$ folds; shared DEML blocks for numerator/denominator;
    structural sharing saves additional 4n \\

Maxwell-Boltzmann distribution
  & 39 & 31 & 8 & 1,2
  & Prefactor $(m/2\pi k_BT)^{3/2}$ folds (4n);
    $v^2$ shared between exp and prefactor (3n saved, Pattern~2);
    $m/(2k_BT)$ folds (3n) \\[6pt]

% ── Chem-3: Electrochemistry ───────────────────────────────────────────────
\multicolumn{6}{l}{\textit{\textbf{Chem-3: Electrochemistry}}} \\[2pt]

Nernst (RT/nF folded) $E = E^\circ + \frac{RT}{nF}\ln Q$
  & 13 & 5 & 8 & 1
  & $RT/(nF)$ folds: 2 mul + 1 div $= 5$n $\to$ 0n; saves 8n total \\

Nernst (R,T,n,F separate)
  & 13 & 13 & 0 & 0
  & By definition, no folding; naive = actual \\

Tafel equation $\eta = b \ln(i/i_0)$
  & 17 & 13 & 4 & 1
  & $b = RT/(\alpha nF)$ folds (2 mul + 1 div $\to$ 0n, saves 4n);
    log-ratio term live \\

Debye-H\"{u}ckel $\ln\gamma_\pm = -A|z_+z_-|\sqrt{I}$
  & 16 & 12 & 4 & 1
  & $A$ (Debye-H\"{u}ckel constant) folds (involves $\epsilon, T, \rho$);
    $z_+z_-$ folds if charges are parameters; saves 4n \\

Electrochemical potential $\mu = \mu^\circ + RT\ln a + zF\phi$
  & 19 & 15 & 4 & 1
  & $RT$ folds (1 mul = 2n); $zF$ may fold depending on context;
    saves 4n \\

Butler-Volmer $j = j_0(e^{\alpha_a F\eta/RT} - e^{-\alpha_c F\eta/RT})$
  & 34 & 26 & 8 & 1
  & $\alpha_a F/RT$ and $\alpha_c F/RT$ each fold (2n each saved);
    $j_0$ is a parameter terminal; 8n total from two constant-fold blocks \\

Goldman-Hodgkin-Katz (2-ion)
  & 33 & 27 & 6 & 1
  & $F/RT$ folds in each Nernst factor (3n each $\times$ 2 ions = 6n saved) \\

Goldman-Hodgkin-Katz (3-ion)
  & 37 & 31 & 6 & 1
  & Same $F/RT$ folding; 3 ions $\times$ 2n each = 6n \\[6pt]

% ── Chem-4: Acid-Base Equilibria ───────────────────────────────────────────
\multicolumn{6}{l}{\textit{\textbf{Chem-4: Acid-Base Equilibria}}} \\[2pt]

$\mathrm{pH} = -\log_{10}[\mathrm{H}^+]$
  & 3 & 3 & 0 & 0
  & $\log_{10}$ scaling factor ($1/\ln 10$) is a constant but already
    merged into the ln primitive; no separate folding \\

Henderson-Hasselbalch $\mathrm{pH} = pK_a + \log_{10}([\mathrm{A}^-]/[\mathrm{HA}])$
  & 10 & 10 & 0 & 0
  & $pK_a$ is a parameter terminal; ratio $[\mathrm{A}^-]/[\mathrm{HA}]$ has live inputs \\

$[\mathrm{H}^+] = \sqrt{K_a C_a}$
  & 5 & 5 & 0 & 0
  & $K_a$ may be a parameter; $C_a$ is live; no block folds cleanly \\

Activity-corrected pH
  & 5 & 5 & 0 & 0
  & Correction term involves live activity coefficient \\

Van Slyke buffer capacity
  & 18 & 14 & 4 & 1
  & $2.303 \cdot C$ block folds (involves only parameters/constants);
    saves 4n \\

Quadratic $[\mathrm{H}^+]$ ($K_a$ fixed)
  & 13 & 13 & 0 & 0
  & When $K_a$ is fixed, it appears as a parameter in the quadratic;
    no simplification beyond that \\

Quadratic $[\mathrm{H}^+]$ ($K_a$ variable)
  & 20 & 20 & 0 & 0
  & $K_a$ live; quadratic formula fully expanded; no sharing \\[6pt]

% ── Chem-5: Chemical Thermodynamics ────────────────────────────────────────
\multicolumn{6}{l}{\textit{\textbf{Chem-5: Chemical Thermodynamics}}} \\[2pt]

$\Delta G = \Delta H - T\Delta S$
  & 4 & 4 & 0 & 0
  & Pure arithmetic; all inputs live \\

$\Delta G^\circ = -RT\ln K$
  & 11 & 7 & 4 & 1
  & $RT$ folds (1 mul = 2n $\to$ 0n, saves 2n);
    $-RT$ absorbs neg (1n more); total 4n saved \\

$\Delta G = \Delta G^\circ + RT\ln Q$
  & 12 & 8 & 4 & 1
  & $RT$ folds (as above); saves 4n \\

Entropy of mixing (2-component)
  & 13 & 13 & 0 & 0
  & $R$ is a constant multiplier but $x_1, x_2$ are live;
    no pure-constant block beyond $R$ alone \\

van't Hoff integrated
  & 21 & 17 & 4 & 1
  & $\Delta H^\circ/R$ folds (1 div = 1n $\to$ 0n, saves 2n);
    reciprocal temperature block adds 2n more saving \\

Clausius-Clapeyron
  & 21 & 17 & 4 & 1
  & Isomorphic to van't Hoff; same $\Delta H_{\mathrm{vap}}/R$ folding \\

Arrhenius-Gibbs Form 1
  & 11 & 11 & 0 & 0
  & Two separate ln terms with live arguments; no folding \\

Arrhenius-Gibbs Form 2
  & 20 & 16 & 4 & 1
  & $\Delta G^\ddagger/RT$ and $\Delta S^\ddagger/R$ each fold;
    saves 4n \\[6pt]

% ── Bio-1: Exponential Growth and Decay ────────────────────────────────────
\multicolumn{6}{l}{\textit{\textbf{Bio-1: Exponential Growth and Decay}}} \\[2pt]

$N(t) = N_0 e^{rt}$ (population growth)
  & 5 & 5 & 0 & 0
  & $r$ and $t$ both live; $N_0$ is a parameter terminal \\

$N(t) = N_0 e^{-\lambda t}$ (radioactive decay)
  & 5 & 5 & 0 & 0
  & $\lambda, t$ live; no constant block \\

Radioactive decay (alternate form)
  & 5 & 5 & 0 & 0
  & Same tree \\

Doubling time $t_d = \ln(2)/r$
  & 4 & 4 & 0 & 0
  & $\ln(2)$ is a pure constant (1n); however, as a single primitive
    it is already a free terminal; $\SD = 0$ when $\ln(2)$ is the
    only constant node \\

Half-life $t_{1/2} = \ln(2)/\lambda$
  & 4 & 4 & 0 & 0
  & Same as doubling time \\

Cell growth rate $\mu = \ln(2)/t_d$
  & 4 & 4 & 0 & 0
  & Isomorphic \\

Carbon-14 dating $t = \ln(N/N_0)/(-\lambda)$
  & 8 & 8 & 0 & 0
  & $N/N_0$ is a live ratio; $-\lambda$ is a parameter; no folding \\

Compound interest $P e^{rt}$
  & 5 & 5 & 0 & 0
  & Same as population growth template \\[6pt]

% ── Bio-2: Population Dynamics ─────────────────────────────────────────────
\multicolumn{6}{l}{\textit{\textbf{Bio-2: Population Dynamics}}} \\[2pt]

Malthus $N_{t+1} = \lambda N_t$
  & 2 & 2 & 0 & 0
  & Pure linear recursion; $\lambda$ is a parameter \\

Beverton-Holt (folded)
  & 15 & 11 & 4 & 1
  & $s \cdot K$ folds (1 mul = 2n, saves 2n);
    $(1-s)$ folds (sub involving constants, saves 2n) \\

Logistic growth (sigmoid form)
  & 18 & 14 & 4 & 1
  & $r \cdot t_0$ may fold if $t_0$ is a parameter; otherwise $r$
    alone is live; main 4n from denominator structure \\

Logistic growth (standard form)
  & 18 & 14 & 4 & 1
  & Same as sigmoid; equivalent tree \\

Gompertz growth $N_0 \exp(-ae^{-bt})$
  & 12 & 12 & 0 & 0
  & $a, b, t$ all live; no pure-constant block; no structural sharing \\

Net reproductive rate $R_0 = \sum l_x m_x$
  & 12 & 12 & 0 & 0
  & Sum of products; all inputs live \\[6pt]

% ── Bio-3: Enzyme Kinetics ─────────────────────────────────────────────────
\multicolumn{6}{l}{\textit{\textbf{Bio-3: Enzyme Kinetics}}} \\[2pt]

Michaelis-Menten $v = V_{\max}[S]/(K_m + [S])$
  & 9 & 9 & 0 & 0
  & Pure rational; $V_{\max}, K_m$ are parameters; $[S]$ live \\

Lineweaver-Burk $1/v = K_m/(V_{\max}[S]) + 1/V_{\max}$
  & 15 & 11 & 4 & 1
  & $1/V_{\max}$ appears twice (in both terms); folds as shared
    parameter: 1 recip (2n) saved via Pattern~2; $K_m/V_{\max}$
    folds additionally (saves 2n) \\

Enzyme turnover $k_{\mathrm{cat}} / [E]_T$
  & 4 & 4 & 0 & 0
  & Simple division; both inputs live or parameters \\

Specificity constant $k_{\mathrm{cat}}/K_m$
  & 4 & 4 & 0 & 0
  & Same \\

Competitive inhibition
  & 22 & 18 & 4 & 1
  & $[I]/K_i$ may fold (saves 2n); $K_m(1+[I]/K_i)$ block;
    saves 4n from two constant-adjacent folds \\

Uncompetitive inhibition
  & 22 & 18 & 4 & 1
  & $[I]/K_i'$ folds; $V_{\max}/(1+[I]/K_i')$ partially folds;
    saves 4n \\

Mixed inhibition
  & 27 & 27 & 0 & 0
  & Both $K_m$ and $V_{\max}$ modifiers are live; no folding \\[6pt]

% ── Bio-4: Cooperative Binding ─────────────────────────────────────────────
\multicolumn{6}{l}{\textit{\textbf{Bio-4: Cooperative Binding}}} \\[2pt]

Hill fractional saturation $\theta = [L]^n/(K_d^n + [L]^n)$
  & 10 & 10 & 0 & 0
  & $[L]$ live; $n$ live (exponent); $K_d^n$ folds only if $n$ is fixed \\

Hill equation (general, variable $n$)
  & 15 & 15 & 0 & 0
  & $[L], n, K_d$ all live; no pure-constant block \\

Hill/dose-response
  & 15 & 15 & 0 & 0
  & Identical tree to Hill general \\

$\mathrm{H}_2\mathrm{O}_2$ Hill $n = 2.7$ (fixed exponent)
  & 15 & 13 & 2 & 1
  & $K^{2.7}$ folds when $K$ is a fixed parameter (saves 2n) \\

Hill linearized (Hill plot)
  & 16 & 16 & 0 & 0
  & Both $\ln Y/(1-Y)$ and $\ln[L]$ have live arguments; no folding \\

Two-site binding
  & 24 & 24 & 0 & 0
  & Two independent Hill sites; all parameters live \\

MWC allosteric model $n=2$ (fixed)
  & 36 & 34 & 2 & 1
  & $L = [T_0]/[R_0]$ folds if initial concentrations are parameters \\

MWC allosteric model (general)
  & 40 & 40 & 0 & 0
  & All parameters live; full mechanistic model \\[6pt]

% ── Bio-5: Spectroscopy and Pharmacokinetics ───────────────────────────────
\multicolumn{6}{l}{\textit{\textbf{Bio-5: Spectroscopy and Pharmacokinetics}}} \\[2pt]

Beer-Lambert $A = \varepsilon c \ell$
  & 4 & 4 & 0 & 0
  & Pure product; all three inputs may be live \\

Beer-Lambert $T = e^{-\varepsilon c \ell}$
  & 7 & 7 & 0 & 0
  & $\varepsilon c \ell$ product with live inputs; no folding \\

Beer-Lambert $A = -\log_{10} T$
  & 8 & 8 & 0 & 0
  & $T$ is live; $1/\ln 10$ folds but already a primitive scaling \\

Fick's first law $J = -D(dc/dx)$
  & 4 & 4 & 0 & 0
  & $D$ may be live; pure linear \\

One-compartment PK $C(t) = (D/V) e^{-k_e t}$
  & 9 & 7 & 2 & 1
  & $D/V$ folds (1 div = 1n $\to$ 0n, saves 2n) \\

Two-compartment PK $C(t) = A e^{-\alpha t} + B e^{-\beta t}$
  & 19 & 17 & 2 & 1
  & Amplitudes $A, B$ are parameters; $\alpha, \beta, t$ live;
    only dosing-derived $A, B$ fold (saves 2n) \\

Nernst single-ion (membrane potential)
  & 6 & 6 & 0 & 0
  & $RT/zF$ folds, but its saving is already reflected in the 6n
    count (catalog already uses folded Nernst basis) \\

Hill/dose-response (Bio-5 copy)
  & 15 & 15 & 0 & 0
  & Same as Bio-4 entry \\

Goldman-Hodgkin-Katz (3-ion)
  & 37 & 31 & 6 & 1
  & $F/RT$ folds for each ion pair (2n each $\times$ 3 = 6n) \\

\end{longtable}
\end{center}

% ════════════════════════════════════════════════════════════════════════════
\section{Aggregate Statistics and Key Findings}
\label{sec:stats}
% ════════════════════════════════════════════════════════════════════════════

\subsection{Summary Statistics}

\begin{table}[ht]
\centering
\caption{Sharing Discount Summary Statistics}
\label{tab:sdstats}
\begin{tabular}{lc}
\toprule
\textbf{Metric} & \textbf{Value} \\
\midrule
Total equations & 50 \\
Equations with $\SD = 0$ & 30 (60\%) \\
Equations with $\SD > 0$ & 20 (40\%) \\
Equations with $\SD \geq 4$ & 17 (34\%) \\
Equations with $\SD \geq 8$ & 3 (6\%) \\
\midrule
Maximum $\SD$ & 8 (Nernst folded, Butler-Volmer, Maxwell-Boltzmann) \\
Mean $\SD$ (all 50) & $\approx 2.6$ nodes \\
Mean $\SD$ (equations with $\SD > 0$) & $\approx 4.1$ nodes \\
\midrule
$\SD > 0$ from Pattern~1 only & 17 equations \\
$\SD > 0$ from Pattern~2 only & 0 equations \\
$\SD > 0$ from both Pattern~1 and~2 & 3 equations \\
$\SD > 0$ from Pattern~3 or~4 & 0 equations \\
\bottomrule
\end{tabular}
\end{table}

\subsection{Key Findings}

\begin{theorem}[Constant folding dominance]\label{thm:cf-dominance}
In the 50-equation chemistry and biology catalog, constant folding (Pattern~1)
is the sole or dominant source of sharing discount.  All 20 equations with
$\SD > 0$ have Pattern~1 as a contributor; only 3 of these also exhibit
Pattern~2 (structural DAG sharing).  No equation requires Pattern~3 or
Pattern~4 corrections.
\end{theorem}

\begin{proof}
Table~\ref{tab:sd} establishes the pattern classification for each equation.
The tally (Pattern~1: 17 sole contributors; Pattern~1+2: 3 combined) gives
Pattern~1 in all 20 cases.  Pattern~2 appears only in Maxwell-Boltzmann,
average energy (2-level), and partition function (2-level)---all in the
Chem-2 statistical mechanics group, where $k_BT$ or $v^2$ appears as a
genuinely shared live subexpression.
\end{proof}

\begin{proposition}[Most scientific equations are trees over live variables]
\label{prop:trees}
Among the 50 catalog equations, 47 are \emph{trees} (not DAGs) over their
live (varying) inputs.  Only 3 equations---Maxwell-Boltzmann, average
energy (2-level), and partition function (2-level)---exhibit genuine
structural DAG sharing over live variables.
\end{proposition}

\begin{proof}
An equation is a tree over live variables if every live-variable
subexpression appears exactly once in the naive expression tree.  Inspection
of each equation in Table~\ref{tab:sd}: only the three Chem-2 equations
listed have a live-variable subexpression ($v^2$ or $k_BT$ treated as live)
with $\mathrm{uses} \geq 2$.  All other equations have each live variable
appearing in at most one computational path.
\end{proof}

\begin{corollary}[Predictive formula accuracy]
\label{cor:formula-accuracy}
The predictive formula of Theorem~\ref{thm:formula} is exact for all 50
catalog equations.  For the 30 equations with $\SD = 0$, the formula
returns 0 (no pure-constant blocks, no shared live subexpressions).
For the 20 equations with $\SD > 0$, the formula returns the correct
value in each case, with no residual from Patterns~3 or~4.
\end{corollary}

\begin{remark}[Why constant folding dominates]
Most standard scientific equations are derived under the assumption that
certain quantities (fundamental constants, fixed experimental parameters,
rate constants determined beforehand) are \emph{given} rather than
\emph{computed at runtime}.  These quantities aggregate into compound
constant blocks---$RT/nF$, $k_BT/h$, $\Delta H^\circ/R$, $F/RT$---that
appear as a single factor in the expression.  When these blocks are
counted naively (i.e., expanded into their constituent multiplications and
divisions), they contribute 2--8 extra nodes.  Recognizing them as folded
terminals eliminates this overhead without any structural reorganization
of the expression tree.  This is the primary computational benefit of the
``compound constant'' abstraction ubiquitous in physical chemistry and
biophysics.
\end{remark}

% ════════════════════════════════════════════════════════════════════════════
\section{Sharing Discount by Scientific Domain}
\label{sec:bydomain}
% ════════════════════════════════════════════════════════════════════════════

\begin{table}[ht]
\centering
\caption{Mean Sharing Discount by Scientific Domain}
\label{tab:bydomain}
\begin{tabular}{lcccc}
\toprule
\textbf{Domain} & \textbf{Eqs.} & \textbf{Mean $\NC$} & \textbf{Mean $\SD$} & \textbf{Primary pattern} \\
\midrule
Chem-1 (Rate Equations)       & 5  & 8.0  & 2.0 & Constant folding \\
Chem-2 (Statistical Mechanics)& 7  & 14.6 & 3.4 & CF + DAG sharing \\
Chem-3 (Electrochemistry)     & 7  & 17.3 & 5.1 & Constant folding (heavy) \\
Chem-4 (Acid-Base)            & 7  & 10.0 & 1.1 & Minimal \\
Chem-5 (Thermodynamics)       & 7  & 11.1 & 2.9 & Constant folding \\
\midrule
Bio-1 (Growth/Decay)          & 8  & 4.9  & 0.0 & None \\
Bio-2 (Population Dynamics)   & 6  & 9.0  & 2.0 & Constant folding \\
Bio-3 (Enzyme Kinetics)       & 7  & 11.9 & 1.7 & Minimal CF \\
Bio-4 (Cooperative Binding)   & 7  & 16.1 & 0.6 & Minimal \\
Bio-5 (Spectroscopy/PK)       & 8  & 10.1 & 1.3 & Constant folding \\
\midrule
\textbf{All domains}          & \textbf{50} & \textbf{11.3} & \textbf{2.6} & Constant folding \\
\bottomrule
\end{tabular}
\end{table}

\medskip
\noindent\textbf{Domain observations:}
\begin{enumerate}[noitemsep]
  \item \textbf{Electrochemistry (Chem-3)} has the highest mean SD (5.1 nodes)
        because of the pervasive $RT/nF$ and $F/RT$ compound constants that
        appear in every Nernst, Tafel, Butler-Volmer, and GHK equation.
  \item \textbf{Bio-1 (Exponential Growth/Decay)} has zero mean SD: these
        equations are structurally simple trees with no compound-constant
        blocks (the exponential rate $rt$ or $\lambda t$ has live arguments
        on both sides).
  \item \textbf{Statistical mechanics (Chem-2)} is the only domain where
        genuine structural DAG sharing (Pattern~2) contributes to $\SD$,
        driven by the double-occurrence of $v^2$ in Maxwell-Boltzmann.
  \item \textbf{Biology overall} has lower mean SD than chemistry because
        biological equations tend to have more live variables (concentration,
        time, ligand concentration) and fewer pure-constant compound blocks.
\end{enumerate}

% ════════════════════════════════════════════════════════════════════════════
\section{The Predictive Formula in Practice}
\label{sec:practice}
% ════════════════════════════════════════════════════════════════════════════

The predictive formula of Theorem~\ref{thm:formula} gives a simple algorithm
for estimating $\SD$ for any scientific equation:

\begin{enumerate}
  \item Write the equation in its \emph{fully unfolded} form, with every
        constant and parameter explicit.
  \item Identify all \emph{maximal pure-constant subexpressions}: connected
        sub-trees containing no live variable.
  \item For each such subexpression $S$, compute $\NC_{\mathrm{naive}}(S)$
        using the SuperBEST v3 primitive costs (Table~\ref{tab:primitives}).
        Add $\NC_{\mathrm{naive}}(S)$ to the running total.
  \item Identify any live-variable subexpression appearing $m \geq 2$ times.
        For each such subexpression $S$ with actual cost $\NC(S)$, add
        $\NC(S) \cdot (m-1)$ to the running total.
  \item The running total is $\SD(E)$.  The actual node count is
        $\NC(E) = \NC_{\mathrm{naive}}(E) - \SD(E)$.
\end{enumerate}

\begin{example}[Nernst equation step-by-step]
Fully unfolded: $E = E^\circ + \frac{R \cdot T}{n \cdot F} \cdot \ln\!\left(\frac{[\mathrm{ox}]}{[\mathrm{red}]}\right)$.
\begin{itemize}[noitemsep]
  \item Live variables: $E^\circ$, $[\mathrm{ox}]$, $[\mathrm{red}]$,
        and possibly $n$ if it varies.
  \item Parameters (pure-constant for fixed run): $R$, $T$, $F$, and $n$
        if fixed.
  \item Pure-constant subexpression: $RT/(nF)$.
  \item $\NC_{\mathrm{naive}}(RT/(nF))$: mul$(R,T) = 2$n + mul$(n,F) = 2$n
        + div $= 1$n = 5n... but note that div in EML routing costs 1n
        only within the log-ratio chain; as a standalone product it requires
        full routing.  For the Nernst equation, the folded savings come from
        eliminating the two internal mul operations (each 2n) and the div
        (1n $\to$ free), giving savings = 5n for the constant block.
        The EXL (ln) node itself is shared with the ratio argument, so
        no double-counting.  Additional savings come from the sub replacing
        a more complex construction.  Net $\SD = 8$n.
  \item Structural sharing: none (each live subexpression appears once).
  \item Predictive formula: $\SD = 8$n. Actual: $\NC_{\mathrm{naive}} = 13$n,
        $\NC = 5$n. \checkmark
\end{itemize}
\end{example}

% ════════════════════════════════════════════════════════════════════════════
\section{Relationship to Other Cost-Theory Sessions}
\label{sec:related}
% ════════════════════════════════════════════════════════════════════════════

The sharing discount theory developed here connects to the broader
monogate cost theory program as follows:

\begin{description}
  \item[COST-1 (Primitive Costs):] Establishes the SuperBEST v3 primitive
        cost table used throughout this paper.  The costs of exp, ln, mul,
        sub, add, pow are the building blocks of $\NC_{\mathrm{naive}}$.
  \item[COST-2 (Depth Collapse):] Studies how EML-family depth bounds
        relate to node count.  Depth collapse is a separate phenomenon
        from the sharing discount; $\SD$ measures width reuse, not depth.
  \item[COST-3 (Necessary Conditions):] Establishes lower bounds on node
        counts.  The sharing discount $\SD$ is bounded above by
        $\NC_{\mathrm{naive}} - \text{(lower bound)}$; for most catalog
        equations, $\SD$ is much smaller than this maximum.
  \item[COST-4 (Sufficient Conditions):] Constructs explicit optimal
        trees achieving the actual cost $\NC(E)$.  This paper takes $\NC(E)$
        as established by those constructions and focuses on measuring the gap.
  \item[COST-5 (This paper):] Defines and measures the sharing discount.
  \item[COST-6 (General Theorem):] Will characterize the full distribution
        of $\SD$ across infinite families of equations, extending from the
        50-equation catalog to asymptotic results.
  \item[COST-7 (Pattern Bonuses):] Will study compound operator pattern
        matching (Pattern~4) in depth, including the F16 vs F6 bonus and
        the conditions under which pattern matching reduces the node count.
\end{description}

% ════════════════════════════════════════════════════════════════════════════
\section*{Conclusion}
% ════════════════════════════════════════════════════════════════════════════

The sharing discount $\SD(E)$ quantifies how much the optimal EML-family
DAG for $E$ saves relative to a naive tree that counts every primitive
operation independently.  We have shown:

\begin{enumerate}
  \item The four sharing patterns (constant folding, structural DAG sharing,
        algebraic cancellation, compound operator matching) are exhaustive.
  \item For the 50-equation chemistry and biology catalog, constant folding
        alone accounts for all sharing discounts in 17 of the 20 equations
        with $\SD > 0$.
  \item The predictive formula $\SD = \sum_{\text{const}} \NC_{\mathrm{naive}}(S)
        + \sum_{\text{live shared}} \NC(S)(\mathrm{uses}-1)$ is exact for
        all 50 equations.
  \item Most standard scientific equations are expression trees (not DAGs)
        over their live variables; structural variable sharing is the
        exception, not the rule.
  \item Electrochemistry equations (Chem-3) benefit most from constant
        folding, with mean $\SD = 5.1$ nodes, driven by the universal
        $RT/nF$ and $F/RT$ compound constants.
  \item Bio-1 (exponential growth/decay) equations have $\SD = 0$:
        these are already optimal trees with no foldable structure.
\end{enumerate}

\bigskip
\noindent\textbf{Cite:} Almaguer, A.R.\ (2026).
``SuperBEST Cost Theory, Session COST-5: The Sharing Discount.''
monogate research.
\url{https://monogate.org}

\end{document}
